\section{The retained spatial bilaplacian and the actual source quotient}
\label{sec:cq-bilaplacian-source}

The original polynomial spatial map and its inverse-fibre trajectory are
\begin{gather*}
 Q(x,y,w)=1+xy,\qquad
 \sigma(x,y,w)=\tfrac12\{Q^3w+y^2Q(4+3xy)\},\\
 \gamma(z)=(z^{-1},-3z/2,13z^2/2),\quad z>0.
\end{gather*}
Here the Euclidean coordinates are exactly $(x,y,w)$, and
$\Delta=\partial_x^2+\partial_y^2+\partial_w^2$.
The symbol $\gamma(z)$ in this section denotes this spatial curve;
the topology $\gamma$ on $\cqE$ retains its earlier definition.
Substitution gives $Q(\gamma(z))=-1/2$ and
$\sigma(\gamma(z))=-z^2/8$.
The pullback
$\sigma^*:\cqO(\cqC)\longrightarrow\cqO(\cqC^3)$,
$h\longmapsto h\circ\sigma$, is a continuous injective algebra
homomorphism for compact convergence. Indeed,
$\eta(s)=(0,0,2s)$ satisfies $\sigma\circ\eta(s)=s$,
so $\eta^*\sigma^*=1$. A compact set has compact image under each
polynomial map, proving continuity of both pullbacks by their
defining supremum seminorms. This also proves that the inverse on
the image is continuous. Restriction to $\cqE$ with its growth
topology is continuous because that topology dominates compact
convergence.

\subsection{Every retained coefficient and the residue functional}
Write $g=\nabla\sigma$, $H=D^2\sigma$, $N=g\cdot g$,
and $L=\Delta\sigma$. For every entire $h$,
\begin{align}
 \Delta(h\circ\sigma)&=Nh''(\sigma)+Lh'(\sigma),\notag\\
 \Delta^2(h\circ\sigma)&=
 N^2h^{(4)}(\sigma)+(4g^{\mathsf T}Hg+2NL)h^{(3)}(\sigma)\notag\\
 &\quad+(2\|H\|_{\rm F}^2+4g\cdot\nabla L+L^2)h''(\sigma)
       +(\Delta L)h'(\sigma).
 \label{eq:cq-bilaplacian-chain}
\end{align}
These formulas use the bilinear polynomial contractions; on the
real locus $\|H\|_{\rm F}^2=\sum_{i,j}H_{ij}^2$ is the usual
squared Frobenius norm. To prove the second identity, the two
product rules give
\begin{align*}
 \Delta(Nh'')&=(\Delta N)h''+
   (2g\cdot\nabla N+NL)h^{(3)}+N^2h^{(4)},\\
 \Delta(Lh')&=(\Delta L)h'+
   (2g\cdot\nabla L+L^2)h''+LN h^{(3)}.
\end{align*}
All derivatives of $h$ are evaluated at $\sigma$. Differentiating
$N=\sum_i(\partial_i\sigma)^2$ gives
$\nabla N=2Hg$ and
$\Delta N=2\sum_{i,j}H_{ij}^2+2g\cdot\nabla L$.
Substitution proves every term in
\eqref{eq:cq-bilaplacian-chain}.

The original derivatives before restriction are
\begin{align*}
 g_1&=\tfrac32yQ^2w+\tfrac12y^3(7+6xy),\\
 g_2&=\tfrac32xQ^2w+\tfrac12(8y+21xy^2+12x^2y^3),\\
 g_3&=\tfrac12Q^3,\\
 L&=3wQ(x^2+y^2)+3y^4+18x^2y^2+21xy+4.
\end{align*}
Differentiating these polynomials and substituting the unchanged
curve yields
\[
 g(\gamma(z))=
 \begin{pmatrix}-9z^3/32\\-3z/16\\-1/16\end{pmatrix},
 \qquad
 H(\gamma(z))=
 \begin{pmatrix}
 -27z^4/4&3z^2/16&-9z/16\\
 3z^2/16&13/4&3/(8z)\\
 -9z/16&3/(8z)&0
 \end{pmatrix}.
\]
The complete contractions are
\begin{align*}
 N(\gamma(z))&=(81z^6+36z^2+4)/1024,\\
 L(\gamma(z))&=(13-27z^4)/4,\\
 \|H(\gamma(z))\|_{\rm F}^2
  &=729z^8/16+9z^4/128+81z^2/128+169/16+9/(32z^2),\\
 (g\cdot\nabla L)(\gamma(z))
  &=9477z^8/512-405z^4/64+27z^2/128+81/32+3/(32z^2),\\
 (g^{\mathsf T}Hg)(\gamma(z))
  &=-2187z^{10}/4096+81z^6/4096
       -81z^4/4096+117z^2/1024+9/1024,\\
 (\Delta L)(\gamma(z))&=-111z^2+36/z^2.
\end{align*}
Consequently the exact restricted operator is
\begin{equation}
 (\Delta^2\sigma^*h)(\gamma(z))
   =\sum_{j=1}^4 C_j(z)h^{(j)}(-z^2/8),
 \label{eq:cq-bilaplacian-all-coefficients}
\end{equation}
where
\begin{align*}
 C_4(z)&=\bigl((81z^6+36z^2+4)/1024\bigr)^2,\\
 C_3(z)&=-6561z^{10}/2048+243z^6/2048
             -135z^4/1024+351z^2/512+31/512,\\
 C_2(z)&=26973z^8/128-4419z^4/64
             +135z^2/64+669/16+15/(16z^2),\\
 C_1(z)&=-111z^2+36/z^2.
\end{align*}
In particular both off-diagonal Hessian occurrences and all
derivatives of $L$ have been retained.

For every entire $h$, continuity of its first four derivatives at
zero and these four Laurent polynomials prove the limit
\begin{equation}
 \mathfrak r(h):=
 \lim_{z\downarrow0}z^2(\Delta^2\sigma^*h)(\gamma(z))
 =36h'(0)+\frac{15}{16}h''(0).
 \label{eq:cq-retained-spatial-residue}
\end{equation}
Indeed $z^2C_4,z^2C_3$ tend to zero, whereas $z^2C_2$ and
$z^2C_1$ tend to $15/16$ and $36$. No evenness assumption on
$h$ is needed. For any $R>0$, Cauchy's integral formula gives
\[
 |\mathfrak r(h)|\le
 \left(\frac{36}{R}+\frac{15}{8R^2}\right)
       \sup_{|s|\le R}|h(s)| .
\]
Thus $\mathfrak r$ is continuous for compact convergence and for
the growth topology. This estimate asserts no continuity in the
unnamed source topology $\tau$.

For the actual coefficient $K_t=T_{\rm D}Z_t$ of
Section~\ref{sec:cq-xi4-heat-returns}, differentiation of its
four original terms gives the exact value
\begin{align}
 \mathfrak r(K_t)
 &=-\frac{59}{275184}
       \left(36Z_t'(3)+\frac{15}{16}Z_t''(3)\right)\notag\\
 &\quad-\frac1{336}
       \left(36Z_t''(3)+\frac{15}{16}Z_t'''(3)\right)\notag\\
 &\quad+\frac1{1296}
       \left(36Z_t'(9/2)+\frac{15}{16}Z_t''(9/2)\right)\notag\\
 &\quad+\frac3{132496}
       \left(36Z_t'(13/2)+\frac{15}{16}Z_t''(13/2)\right).
 \label{eq:cq-xi4-spatial-residue}
\end{align}
For the unchanged incidence factor $a_{ef}$ and physical
diffusion $\nu\Delta$, its scaled second diffusion is
$\nu^2a_{ef}\mathfrak r(K_t)$ by linearity and operator composition.
The limit probes spatial infinity: $|\gamma(z)|\ge z^{-1}$.
Every $\sigma^*h$ is smooth at every finite real point by its
entire composition. These assertions specify the domain and
meaning of the observable exactly.

\subsection{The full ambiguity on the actual source quotient}
Set $E=\cqE$, $N=\cqN$, and $Q_{\rm ar}=(E,\tau)/N$.
The original even function $\Xi$ has $\Xi(0)>0$, as proved
from its retained positive kernel. Define
\begin{equation}
 e(s)=\frac{8}{15\Xi(0)}s^2\Xi(s),\qquad
 K=\ker\mathfrak r,\qquad P f=f-e\,\mathfrak r(f).
 \label{eq:cq-residue-section}
\end{equation}
Then $e\in\cqC[s]\Xi\subset N$, $e'(0)=0$, and
$e''(0)=16/15$, so $\mathfrak r(e)=1$.
Consequently
\begin{equation}
 \mathfrak r(N)=\cqC,\qquad
 \{([f]_N,\mathfrak r(f)):f\in E\}
       =Q_{\rm ar}\times\cqC .
 \label{eq:cq-source-residue-full-ambiguity}
\end{equation}
Here is the exact representative for every pair on the right.
Given any representative $f$ of $q$ and any $a\in\cqC$, put
$f_a=f+(a-\mathfrak r(f))e$. Then $[f_a]_N=q$ and
$\mathfrak r(f_a)=a$. In particular the zero class itself has
representatives $ae$ with every residue $a$. No nonempty set of
source quotient classes admits this scalar as a value independent
of all its representatives.

The strongest corresponding quotient map is explicit. Algebraically
$P^2=P$, $\operatorname{im}P=K$, $\ker P=\cqC e$, and
\[
 N_0:=N\cap K=P(N),\qquad P^{-1}(N_0)=N .
\]
The first equality for $N_0$ follows because $e\in N$ and
$P$ fixes $K$. If $Pf\in N$, the identity
$f=Pf+e\mathfrak r(f)$ proves $f\in N$, establishing the
inverse-image equality. Give $K$ the quotient topology
$\tau_P=P_*\tau$: a subset $V\subset K$ is open exactly when
$P^{-1}(V)$ is open in $(E,\tau)$. This is the topology of
$(E,\tau)/\cqC e$ under its displayed algebraic isomorphism
with $K$. In particular it is a vector topology. Since $N$ is
closed, the same inverse-image equality proves that $N_0$ is
closed in $(K,\tau_P)$. There is a homeomorphism
\begin{equation}
 (K,\tau_P)/N_0
    \ \xrightarrow{\ \ [k]\mapsto[k]_N\ \ }\
    (E,\tau)/N .
 \label{eq:cq-source-residue-reduced-quotient}
\end{equation}
For completeness, let $q_N:E\to E/N$ and
$q_0:K\to K/N_0$ denote the quotient maps. The proposed map
$A$ satisfies $Aq_0P=q_N$. Since $q_0P$ is a composition
of quotient maps and hence is quotient, this identity proves
continuity of $A$. The inverse is induced by $q_0P$, because
its kernel is $N$; the quotient property of $q_N$ proves its
continuity. These maps are inverse since
$P f-f\in N$ and $P k=k$. This proves the assertion without
using $\tau$-continuity of $\mathfrak r$ or of $P$ as an
endomorphism of $(E,\tau)$.

There is also a proved comparison using the inherited topology:
the identity $(K,\tau|_K)\to(K,\tau_P)$ is continuous.
Indeed for $V$ open in $\tau_P$,
$V=K\cap P^{-1}(V)$, which is inherited-open.
It induces a continuous bijection
$(K,\tau|_K)/N_0\to Q_{\rm ar}$.
No equality of these two topologies on $K$ is asserted.
For the growth topology the stronger result is available:
$\mathfrak r$ and $P$ are continuous by the Cauchy bound,
so $f\mapsto(Pf,\mathfrak r(f))$ and
$(k,a)\mapsto k+ae$ are inverse continuous linear maps
\begin{equation}
 (E,\gamma)\simeq (K,\gamma|_K)\oplus\cqC .
 \label{eq:cq-residue-growth-splitting}
\end{equation}
The continuity of $a\mapsto ae$ follows from each defining
seminorm. This proves the full split topology at precisely
the topology where its hypotheses have been established.

\subsection{Every finite jet at the original coordinate origin}
The residue calculation is a particular finite-jet phenomenon.
For $d\ge0$ retain the derivative coordinates
$J_d f=(f(0),f'(0),\ldots,f^{(d)}(0))$.
Write the convergent germ
$1/\Xi(s)=\sum_{n\ge0}b_ns^n$; its coefficients are given
without any change of $\Xi$ by
\[
 b_0=\Xi(0)^{-1},\qquad
 b_n=-\Xi(0)^{-1}
       \sum_{k=1}^n\frac{\Xi^{(k)}(0)}{k!}\,b_{n-k}.
\]
For $a=(a_0,\ldots,a_d)$ define the entire polynomial trace
\begin{equation}
 E_da(s)=\Xi(s)\sum_{n=0}^d
           \left(\sum_{j=0}^n\frac{a_j}{j!}b_{n-j}\right)s^n.
 \label{eq:cq-origin-jet-section}
\end{equation}
Multiplying its truncated Taylor series by that of $\Xi$
gives $E_da(s)=\sum_{j=0}^d a_js^j/j!+O(s^{d+1})$,
because the defining reciprocal convolution has coefficient
one at zero and zero at each positive degree.
Therefore $J_dE_d=1$ and
$\operatorname{im}E_d\subset\cqC[s]\Xi\subset N$.
This proves $J_d(N)=\cqC^{d+1}$, with its complete
finite-dimensional section, rather than only the scalar image.
For any $q=[f]_N$ and any jet $a$, the representative
$f+E_d(a-J_df)$ realizes $(q,a)$.

Put $P_d=1-E_dJ_d$, $K_d=\ker J_d$, and
$N_d=N\cap K_d$. The same identities, now with
$\ker P_d=\operatorname{im}E_d$, prove
$P_d(N)=N_d$ and $P_d^{-1}(N_d)=N$.
Thus $(K_d,(P_d)_*\tau)/N_d\simeq Q_{\rm ar}$ by the
two quotient-map arguments just given. Cauchy's bounds
$|f^{(j)}(0)|\le j!R^{-j}\sup_{|s|\le R}|f(s)|$
prove the continuous growth-topology decomposition
$(E,\gamma)\simeq (K_d,\gamma|_{K_d})\oplus\cqC^{d+1}$.
All factorials are retained in the derivative coordinates
and in the section.

\subsection{Heat transport of the retained residue and its quotient}
The scalar observable also has an exact heat transport. Define
\[
 \mathfrak r_t=\mathfrak r U_{-t},\qquad
 e_t=U_te,\qquad
 P_t=1-e_t\mathfrak r_t=U_tPU_{-t},\qquad
 K_t^{\rm res}=\ker\mathfrak r_t .
\]
These definitions retain $U_t=\exp(-tD^2/4)$ and hence
give the convergent formula
\begin{equation}
 \mathfrak r_t(f)
 =36\sum_{k=0}^{\infty}\frac{(t/4)^k}{k!}
       f^{(2k+1)}(0)
  +\frac{15}{16}\sum_{k=0}^{\infty}\frac{(t/4)^k}{k!}
       f^{(2k+2)}(0).
 \label{eq:cq-transported-spatial-residue}
\end{equation}
The earlier heat-series proof gives convergence in the growth
topology, locally uniformly in $t$; applying the two continuous
derivative evaluations gives the displayed convergence and
its value. In particular
$\mathfrak r_tU_t=\mathfrak r$,
$\mathfrak r_t(e_t)=1$, and $e_t\in N_t$.
Also $U_tK=K_t^{\rm res}$ and
$U_tN_0=N_t\cap K_t^{\rm res}$ by these identities and
invertibility of $U_t$.

Give $K_t^{\rm res}$ the topology $(P_t)_*\tau_t$.
It equals $(U_t)_*\tau_P$: for $V\subset K_t^{\rm res}$,
the identity $P_tU_t=U_tP$ gives
$U_{-t}(P_t^{-1}(V))=P^{-1}(U_{-t}V)$,
which is exactly the equivalence of the two defining open-set
conditions. Thus $U_t$ gives a homeomorphism between the
residue-zero quotient in
\eqref{eq:cq-source-residue-reduced-quotient} and
\[
 (K_t^{\rm res},(P_t)_*\tau_t)/
       (N_t\cap K_t^{\rm res}) .
\]
The latter is homeomorphic to $(E,\tau_t)/N_t$ by
$[k]\mapsto[k]_{N_t}$, with inverse $[f]\mapsto[P_tf]$;
the two quotient universal properties apply exactly as above.
The two homeomorphisms commute with the already proved heat
homeomorphism of the actual source quotient because both send
$[f]$ to $[U_tf]$. Finally, for each $q=[f]_{N_t}$ and
$a\in\cqC$, the unchanged class has the representative
$f+e_t(a-\mathfrak r_t(f))$ with transported residue $a$.
This computes the full transported scalar relation without
identifying $\mathfrak r_t$ with the fixed physical functional
$\mathfrak r$.
